\documentclass[10pt, twocolumn]{article}

% ── Packages ────────────────────────────────────────────────────────────────
\usepackage[T1]{fontenc}
\usepackage[utf8]{inputenc}
\usepackage{lmodern}
\usepackage[a4paper, top=2.2cm, bottom=2.2cm, left=1.6cm, right=1.6cm,
            columnsep=0.6cm]{geometry}
\usepackage{amsmath, amssymb}
\usepackage{graphicx}
\usepackage{booktabs}
\usepackage{array}
\usepackage{xcolor}
\usepackage{hyperref}
\usepackage{enumitem}
\usepackage{microtype}
\usepackage{caption}
\usepackage{float}
\usepackage{authblk}
\usepackage{titlesec}
\usepackage{fancyhdr}
\usepackage{abstract}
\usepackage{parskip}
\usepackage{listings}
\usepackage{cite}

% ── Colours ──────────────────────────────────────────────────────────────────
\definecolor{aims-blue}{RGB}{0, 70, 127}
\definecolor{code-bg}{RGB}{245, 245, 245}

% ── Hyperref ─────────────────────────────────────────────────────────────────
\hypersetup{
  colorlinks = true,
  linkcolor  = aims-blue,
  urlcolor   = aims-blue,
  citecolor  = aims-blue,
}

% ── Section title formatting ─────────────────────────────────────────────────
\titleformat{\section}{\normalfont\large\bfseries\color{aims-blue}}
  {\thesection}{0.6em}{}
\titleformat{\subsection}{\normalfont\normalsize\bfseries}
  {\thesubsection}{0.5em}{}
\titleformat{\subsubsection}{\normalfont\small\bfseries\itshape}
  {\thesubsubsection}{0.5em}{}

% ── Header / footer ──────────────────────────────────────────────────────────
\pagestyle{fancy}
\fancyhf{}
\lhead{\small\color{gray} AIMS RIC — Doctoral Training School 2026}
\rhead{\small\color{gray} ML-ESS Technical Report}
\cfoot{\small\thepage}
\renewcommand{\headrulewidth}{0.4pt}

% ── Abstract formatting ──────────────────────────────────────────────────────
\renewcommand{\abstractnamefont}{\normalfont\bfseries}
\renewcommand{\abstracttextfont}{\normalfont\small}
\setlength{\absleftindent}{0pt}
\setlength{\absrightindent}{0pt}

% ── Caption formatting ───────────────────────────────────────────────────────
\captionsetup{font=small, labelfont=bf, skip=4pt}

% ── Listings ─────────────────────────────────────────────────────────────────
\lstset{
  basicstyle=\ttfamily\footnotesize,
  backgroundcolor=\color{code-bg},
  frame=single, framerule=0.4pt,
  breaklines=true, breakatwhitespace=true,
  captionpos=b, xleftmargin=4pt, xrightmargin=4pt,
}

% ── List spacing ─────────────────────────────────────────────────────────────
\setlist[itemize]{noitemsep, topsep=2pt, leftmargin=*}
\setlist[enumerate]{noitemsep, topsep=2pt, leftmargin=*}

% ─────────────────────────────────────────────────────────────────────────────
%  TITLE BLOCK
% ─────────────────────────────────────────────────────────────────────────────

\title{%
  \vspace{-1em}%
  {\large\textbf{AIMS RIC — Doctoral Training School 2026}}\\[4pt]
  {\Large\bfseries\color{aims-blue}
    ML-ESS: A Multi-Agent AI Research Assistant}\\[4pt]
  {\normalsize Technical Report — April 2026}%
}

\author[1]{\textbf{Jérémie N. Landu}}
\affil[1]{%
  African Institute for Mathematical Sciences\\
  Research Innovation Centre (AIMS-RIC)\\
  \texttt{jeremy@aimsammi.org}
}

\date{}

% =============================================================================
\begin{document}
% =============================================================================

\maketitle
\thispagestyle{fancy}

% ─────────────────────────────────────────────────────────────────────────────
\begin{abstract}
ML-ESS is a full-stack AI system that automates end-to-end research: given a
natural language question, it orchestrates four specialised AI agents to search
the web, synthesise evidence, draft a structured report, and evaluate the
report's quality. The system exposes a REST API (FastAPI) and a modern web
interface (Next.js), supporting real-time progress tracking, PDF export, and
multi-provider LLM fallback. This report describes the system architecture,
agent pipeline, data model, REST API, frontend application, deployment
strategy, and design trade-offs.
\end{abstract}

\vspace{0.5em}
\noindent\textbf{Keywords:} multi-agent systems, large language models,
automated research, FastAPI, Next.js, retrieval-augmented generation.

\vspace{0.5em}
\hrule
\vspace{0.5em}

% =============================================================================
\section{Introduction \& Motivation}
% =============================================================================

Conducting thorough research manually is time-consuming: it involves querying
multiple sources, extracting relevant claims, resolving contradictions, and
composing a coherent, well-cited document. ML-ESS automates this entire
workflow through a multi-agent pipeline, replacing hours of manual work with a
structured, reproducible process.

\medskip
\noindent\textbf{Core objectives:}
\begin{itemize}
  \item Accept any research question in natural language.
  \item Autonomously gather evidence from the open web.
  \item Synthesise findings into themes, resolving contradictions.
  \item Generate a structured Markdown report with citations.
  \item Score the report quality along four dimensions.
  \item Deliver everything through a clean web UI with real-time feedback.
\end{itemize}

% =============================================================================
\section{System Architecture}
% =============================================================================

\subsection{High-Level Overview}

The system is divided into two main services: a Python backend responsible for
all AI reasoning and data persistence, and a Next.js frontend that provides the
user interface. Figure~\ref{fig:arch} illustrates the high-level architecture.

\begin{figure}[H]
  \centering
  \includegraphics[width=\linewidth]{../report-rendered-1.png}
  \caption{High-level system architecture.}
  \label{fig:arch}
\end{figure}

\subsection{Backend Directory Structure}

\begin{lstlisting}
api/
|-- main.py                 # CLI & server entry point
|-- app/
    |-- agents/
    |   |-- search.py       # Web search & evidence extraction
    |   |-- synthesis.py    # Theme identification
    |   |-- report.py       # Markdown generation
    |   |-- evaluator.py    # Quality scoring
    |-- api/
    |   |-- app.py          # FastAPI factory, CORS
    |   |-- routes.py       # REST endpoints & SSE
    |   |-- auth.py         # API key authentication
    |-- core/
    |   |-- pipeline.py     # Sequential orchestration
    |   |-- jobs.py         # Background thread management
    |   |-- store.py        # SQLite (WAL mode)
    |   |-- llm.py          # LLM client w/ retries
    |-- models/
        |-- api.py          # Pydantic schemas
        |-- state.py        # Pipeline shared state
\end{lstlisting}

\subsection{Request Lifecycle \& Data Flow}

Figure~\ref{fig:lifecycle} shows the full request lifecycle from the moment a
user submits a question to the delivery of the final report.

\begin{figure}[H]
  \centering
  \includegraphics[width=\linewidth]{../report-rendered-2.png}
  \caption{Request lifecycle and data flow (sequence diagram).}
  \label{fig:lifecycle}
\end{figure}

% =============================================================================
\section{The Multi-Agent Pipeline}
% =============================================================================

Each agent in the pipeline is a stateless function that reads from and writes
to a shared \texttt{SharedState} object. Agents are called sequentially — each
one enriches the state for the next.

\subsection{SharedState Schema}

Figure~\ref{fig:state} shows the class diagram of the shared pipeline state,
including the \texttt{Source}, \texttt{Evidence}, \texttt{Theme},
\texttt{Contradiction}, and \texttt{EvaluationScores} entities.

\begin{figure}[H]
  \centering
  \includegraphics[width=\linewidth]{../report-rendered-3.png}
  \caption{SharedState class diagram.}
  \label{fig:state}
\end{figure}

\subsection{Agent Pipeline Flow}

The pipeline executes four agents sequentially
(Figure~\ref{fig:pipeline}):

\begin{enumerate}
  \item \textbf{Search Agent} — generates queries, fetches pages, and extracts
    structured evidence.
  \item \textbf{Synthesis Agent} — groups evidence into themes and detects
    contradictions.
  \item \textbf{Report Agent} — produces a structured Markdown report
    (\textasciitilde{}2\,000–2\,500 words) with citations.
  \item \textbf{Evaluator} — scores the report on four quality dimensions.
\end{enumerate}

\begin{figure}[H]
  \centering
  \includegraphics[width=\linewidth]{../report-rendered-4.png}
  \caption{Agent pipeline flow.}
  \label{fig:pipeline}
\end{figure}

\subsection{Search Agent Detail}

The Search Agent (Figure~\ref{fig:search}) generates up to three search
queries via an LLM call, fans out to DuckDuckGo (or Tavily when an API key is
provided), de-duplicates the result URLs, fetches page content with
\texttt{httpx}, parses it with BeautifulSoup4, and finally calls the LLM again
to extract structured evidence records.

\begin{figure}[H]
  \centering
  \includegraphics[width=\linewidth]{../report-rendered-5.png}
  \caption{Search Agent internal data flow.}
  \label{fig:search}
\end{figure}

\subsection{LLM Provider Strategy}

Figure~\ref{fig:llm} shows the retry and fallback strategy implemented in the
LLM client. The primary provider is Groq (\texttt{llama-3.3-70b-versatile});
HuggingFace Inference (\texttt{Llama-3.3-70B-Instruct}) serves as a fallback
when Groq rate-limits are exhausted.

\begin{figure}[H]
  \centering
  \includegraphics[width=\linewidth]{../report-rendered-6.png}
  \caption{LLM provider retry and fallback strategy.}
  \label{fig:llm}
\end{figure}

\subsection{Evaluation Scores}

The Evaluator Agent scores the generated report along four dimensions
(Table~\ref{tab:scores}). Scores for coverage, faithfulness, and usefulness
are in $[0,1]$; the hallucination rate is also in $[0,1]$ but lower is better.

\begin{table}[H]
  \centering
  \small
  \caption{Example evaluation scores.}
  \label{tab:scores}
  \begin{tabular}{lc}
    \toprule
    \textbf{Dimension} & \textbf{Score} \\
    \midrule
    Coverage          & 0.87 \\
    Faithfulness      & 0.92 \\
    Usefulness        & 0.85 \\
    Hallucination rate & 0.08 \\
    \bottomrule
  \end{tabular}
\end{table}

% =============================================================================
\section{Data Persistence}
% =============================================================================

\subsection{SQLite Schema}

The backend persists all job data in a single SQLite database opened in WAL
mode (Figure~\ref{fig:er}). Three tables are used: \texttt{JOBS} (primary
record), \texttt{STATES} (full serialised pipeline state), and \texttt{EVENTS}
(append-only SSE event log).

\begin{figure}[H]
  \centering
  \includegraphics[width=\linewidth]{../report-rendered-7.png}
  \caption{Entity-relationship diagram.}
  \label{fig:er}
\end{figure}

\subsection{Job Status Lifecycle}

Figure~\ref{fig:states} shows the finite-state machine for a research job. A
job progresses through \textit{pending} $\to$ \textit{searching} $\to$
\textit{synthesising} $\to$ \textit{reporting} $\to$ \textit{evaluating} $\to$
\textit{completed}, or transitions to \textit{failed} on any unhandled
exception.

\begin{figure}[H]
  \centering
  \includegraphics[width=\linewidth]{../report-rendered-8.png}
  \caption{Job status state machine.}
  \label{fig:states}
\end{figure}

\subsection{SSE Event Stream}

Figure~\ref{fig:sse} shows the Server-Sent Events (SSE) protocol between the
frontend and the \texttt{/api/research/\{id\}/events} endpoint.

\begin{figure}[H]
  \centering
  \includegraphics[width=\linewidth]{../report-rendered-9.png}
  \caption{SSE event stream sequence.}
  \label{fig:sse}
\end{figure}

% =============================================================================
\section{REST API Reference}
% =============================================================================

All endpoints except \texttt{/api/health} require the header
\texttt{X-API-Key: <key>} when \texttt{API\_KEY} is configured.

\begin{table}[H]
  \centering
  \small
  \caption{REST API endpoints.}
  \label{tab:api}
  \begin{tabular}{p{1.0cm} p{3.6cm} c}
    \toprule
    \textbf{Method} & \textbf{Endpoint} & \textbf{Status} \\
    \midrule
    \texttt{POST}   & \texttt{/api/research}            & 202 \\
    \texttt{GET}    & \texttt{/api/research}             & 200 \\
    \texttt{GET}    & \texttt{/api/research/\{id\}}      & 200 \\
    \texttt{GET}    & \texttt{.../\{id\}/reasoning}      & 200 \\
    \texttt{GET}    & \texttt{.../\{id\}/events}         & 200 \\
    \texttt{GET}    & \texttt{.../\{id\}/pdf}            & 200 \\
    \texttt{DELETE} & \texttt{/api/research/\{id\}}      & 200 \\
    \texttt{DELETE} & \texttt{/api/research}             & 200 \\
    \texttt{GET}    & \texttt{/api/health}               & 200 \\
    \bottomrule
  \end{tabular}
\end{table}

% =============================================================================
\section{Frontend Application}
% =============================================================================

The Next.js~16 frontend (React~19, Tailwind CSS~v4, TypeScript) provides three
pages: Home (\texttt{/}), Job History (\texttt{/history}), and Job Detail
(\texttt{/job/[id]}). Figure~\ref{fig:fe-tree} shows the page and component
tree; Figure~\ref{fig:fe-flow} shows the frontend data flow.

\begin{figure}[H]
  \centering
  \includegraphics[width=\linewidth]{../report-rendered-10.png}
  \caption{Frontend page and component tree.}
  \label{fig:fe-tree}
\end{figure}

\begin{figure}[H]
  \centering
  \includegraphics[width=\linewidth]{../report-rendered-11.png}
  \caption{Frontend data flow.}
  \label{fig:fe-flow}
\end{figure}

% =============================================================================
\section{Technology Stack}
% =============================================================================

\begin{table}[H]
  \centering
  \small
  \caption{Technology stack summary.}
  \label{tab:stack}
  \begin{tabular}{p{1.8cm} p{2.0cm} p{2.2cm}}
    \toprule
    \textbf{Layer} & \textbf{Technology} & \textbf{Role} \\
    \midrule
    Backend lang.    & Python 3.11+      & Agent \& API logic \\
    API framework    & FastAPI $\geq$0.115 & HTTP, routing \\
    ASGI server      & Uvicorn $\geq$0.30 & Production server \\
    LLM (primary)    & Groq SDK          & llama-3.3-70b \\
    LLM (fallback)   & HuggingFace Hub   & Llama-3.3-70B \\
    Web search       & ddgs $\geq$9.0    & DuckDuckGo \\
    Web search (opt) & tavily-python     & Premium search \\
    Scraping         & httpx + bs4       & Fetch \& parse \\
    Validation       & Pydantic $\geq$2  & Schemas \\
    Database         & SQLite (WAL)      & Persistence \\
    PDF export       & WeasyPrint $\geq$62 & MD $\to$ PDF \\
    Frontend         & Next.js 16.2.1    & React 19 App Router \\
    Styling          & Tailwind CSS v4   & Utility-first CSS \\
    Language (FE)    & TypeScript        & Type safety \\
    Deployment (FE)  & Vercel            & Hosted frontend \\
    \bottomrule
  \end{tabular}
\end{table}

% =============================================================================
\section{Configuration}
% =============================================================================

Key environment variables (set in \texttt{api/.env}) are listed in
Table~\ref{tab:config}.

\begin{table}[H]
  \centering
  \small
  \caption{Key environment variables.}
  \label{tab:config}
  \begin{tabular}{p{2.6cm} p{1.1cm} p{2.5cm}}
    \toprule
    \textbf{Variable} & \textbf{Default} & \textbf{Purpose} \\
    \midrule
    \texttt{LLM\_PROVIDER}         & auto      & groq / huggingface / auto \\
    \texttt{GROQ\_API\_KEY}        & —         & Primary Groq key \\
    \texttt{GROQ\_API\_KEY\_2}     & —         & Fallback Groq key \\
    \texttt{GROQ\_MODEL}           & llama-3.3-70b & Groq model \\
    \texttt{HF\_TOKEN}             & —         & HuggingFace token \\
    \texttt{TEMPERATURE}           & 0.3       & LLM temperature \\
    \texttt{TAVILY\_API\_KEY}      & —         & Premium search \\
    \texttt{MAX\_SEARCH\_QUERIES}  & 3         & Queries per question \\
    \texttt{MAX\_RESULTS\_PER\_QUERY} & 5      & Results per query \\
    \texttt{MAX\_CONTENT\_LENGTH}  & 3000      & Chars per page \\
    \texttt{API\_KEY}              & —         & Auth key \\
    \texttt{DB\_PATH}              & data/jobs.db & SQLite location \\
    \texttt{CORS\_ORIGINS}         & Vercel URL & Allowed origins \\
    \bottomrule
  \end{tabular}
\end{table}

% =============================================================================
\section{Deployment Architecture}
% =============================================================================

\subsection{Current vs.\ Scalable Architecture}

Figure~\ref{fig:deploy} contrasts the current single-instance deployment
(FastAPI + SQLite + background threads on a single process) with a
production-scale horizontal architecture (load balancer, multiple FastAPI
workers, Redis job queue, Celery workers, and PostgreSQL).

\begin{figure}[H]
  \centering
  \includegraphics[width=\linewidth]{../report-rendered-12.png}
  \caption{Current vs.\ scalable deployment architecture.}
  \label{fig:deploy}
\end{figure}

\subsection{Deployment Checklist}

Figure~\ref{fig:checklist} summarises the five-step deployment procedure:
set environment variables, initialise SQLite, start Uvicorn, configure the
frontend, and verify the health endpoint.

\begin{figure}[H]
  \centering
  \includegraphics[width=\linewidth]{../report-rendered-13.png}
  \caption{Deployment checklist flowchart.}
  \label{fig:checklist}
\end{figure}

% =============================================================================
\section{Conclusion}
% =============================================================================

ML-ESS demonstrates how a small, well-structured multi-agent system can
automate a complex, knowledge-intensive task end-to-end. By composing four
focused agents — each with a clear role — the pipeline transforms a single
natural language question into a cited, evaluated research report in minutes.

The system is designed for clarity and correctness: low LLM temperature for
deterministic outputs, retry/fallback logic for resilience, persistent job
state for observability, and a clean typed API for frontend integration. The
modular agent design makes the system straightforward to extend: adding a new
research phase (e.g., a fact-checking agent or a translation layer) requires
only inserting a new agent function into the pipeline and extending
\texttt{SharedState} with the new fields.

% =============================================================================
%  ACKNOWLEDGEMENTS
% =============================================================================

\section*{Acknowledgements}

The author thanks the AIMS Research Innovation Centre and the organisers of the
2026 Doctoral Training School for providing the research environment and
computational resources that made this work possible.

% =============================================================================
%  REFERENCES  (add a .bib file or use manual \bibitem entries)
% =============================================================================

\begin{thebibliography}{9}

\bibitem{fastapi}
S. Ramirez, \textit{FastAPI}, 2019. \url{https://fastapi.tiangolo.com}

\bibitem{nextjs}
Vercel, \textit{Next.js Documentation}, 2024. \url{https://nextjs.org/docs}

\bibitem{groq}
Groq Inc., \textit{Groq API — llama-3.3-70b-versatile}, 2024.
\url{https://console.groq.com}

\bibitem{llama3}
AI Meta, ``The Llama 3 Herd of Models,'' \textit{arXiv preprint
arXiv:2407.21783}, 2024.

\bibitem{rag}
P. Lewis \textit{et al.}, ``Retrieval-Augmented Generation for
Knowledge-Intensive NLP Tasks,'' in \textit{Advances in Neural Information
Processing Systems}, vol. 33, 2020, pp. 9459--9474.

\bibitem{weasyprint}
CourtBouillon, \textit{WeasyPrint}, 2024. \url{https://weasyprint.org}

\end{thebibliography}

\end{document}
